% Avoid simply summarizing the findings of each paper, one paragraph per paper. \cite{hosseini2016adaptive} discusses.. \cite{SURV} writes  

% Instead, synthesize the papers by identifying themes and using the papers as references for the themes. 

% NEED TO have conference co chairs if citing a conference paper

% journals from the last 2-3 years
The architectural approach of this research is heavily informed by previous exploration into hardware-software co-design for implantable medical devices. Historically, the fundamental challenge of respecting the strict $1^\circ$C thermal limit has forced designers to rely on highly specialized, monolithic ASICs. The NeuroPace RNS System \cite{Heck2014}, Neuralink's N1 ASIC \cite{Musk2019}, and Kassiri et al.'s closed-loop neurostimulator \cite{7982670} represent the state of the art implementations of this rigid baseline. These systems achieve unparalleled energy efficiency for their exact medical or recording tasks by permanently baking clinical logic into the physical silicon. However, they inherently lack flexibility in case a neuroscientist wants to update any of the logic or algorithms. This project directly contrasts this rigid historical approach by employing a software-defined, architecture-specific compilation pipeline built on the HALO platform. As demonstrated in the Results section, the flexible RISC-V software pipeline allows the system to instantly benefit from ongoing toolchain enhancements without any hardware modifications. Simply updating the GCC cross-compiler yielded a collective execution improvement of 7.2 percent across the core arithmetic and logical operations (excluding Octave-mode comparison tests). Furthermore, the pipeline gracefully handles complex, dynamic software logic—such as the early-exit short-circuiting behavior discovered in the isEqual complex matrix tests. If a neuroscientist needed to optimize a comparison algorithm or fix a logic bug on the monolithic ASICs cited above, a completely new tape-out would be required; on the HALO RISC-V pipeline, it simply requires a software recompile.

\section{Source-to-Source Translation for High-Level Languages}

In the domain of high-level synthesis and automated code generation, commercial tools like MathWorks' MATLAB Coder \cite{mathworks_matlab_coder} provide a highly refined reference point. MATLAB Coder is highly capable in parsing proprietary MATLAB scripts and generating optimized C and C++ code targeted primarily at advanced desktop processors and commercial ARM architectures. Similarly, for researchers utilizing Python—another dynamically typed language widely used in BCI data analysis and real-time processing pipelines—tools like Cython \cite{behnel2011cython} offer similar benefits. Cython addresses the performance bottlenecks of vanilla Python by compiling code directly into efficient C or C++ extensions, utilizing static type definitions to achieve substantial execution speedups for data pipelines.

While both MATLAB Coder and Cython are highly effective for accelerating signal processing on standard architectures, they share a critical limitation in the context of this research: they are structurally agnostic regarding emerging, open-source instruction sets. Neither tool inherently supports the automated generation of targeted RISC-V Vector (RVV) intrinsics required to meet the ultra-low-power constraints of the HALO platform. Consequently, the custom ts-traversal and RISCV-Matrix integration developed in this work is uniquely necessary to bridge the semantic gap between high-level BCI algorithm prototyping and the specific vector execution demands of the embedded hardware.
